\documentclass[11pt,a4paper]{article}
\usepackage[T1]{fontenc}
\usepackage{lmodern,microtype,amsmath,amssymb,amsthm,booktabs}
\usepackage[hmargin=26mm,vmargin=22mm]{geometry}
\usepackage[hidelinks]{hyperref}
\newtheorem{theorem}{Theorem}
\newtheorem*{corollary}{Corollary}
\newcommand{\F}{\mathbb F}
\newcommand{\N}{\mathbb N}
\title{What stopping-time counts can decide\\[2pt]\large A convergence criterion, a divergent polynomial orbit, and an integer basin}
\author{}\date{}
\begin{document}
\maketitle\vspace{-3em}
\section{A weaker target for an open integer system}
Let $T$ on the positive integers be the census map \texttt{split5\_3},
\[
 T(n)=\begin{cases}n/2&2\mid n,\\5n+3& n\equiv1\pmod4,\\n+3&n\equiv3\pmod4.\end{cases}
\]
Write $N(X,b)$ for the number of $1\le n\le X$ surviving the first $b$ halvings without a strict drop below $n$, using actual orbit heights. The integer count conjecture E predicts this quantity up to small relative error and a finite allowance for permanent survivors. To prove that all orbits eventually cycle, even a logarithmically growing allowance can suffice.
\begin{theorem}\label{thm:records}
If some positive orbit of $T$ never becomes periodic, then for every sufficiently large $X$ and every depth $b$,
\[
  N(X,b)\ \ge\ \frac{\log_2 X}{\log_2(5/4)}-O(1).
\]
The additive constant can depend on the orbit. In particular, for any chosen depths $b_m$, the single inequality
\[
 \liminf_{m\to\infty}\frac{N(2^m,b_m)}m
 < \frac1{\log_2(5/4)}=3.10628\ldots
\]
implies that every positive orbit of $T$ eventually cycles.
\end{theorem}
\begin{proof}
An orbit that never becomes periodic never repeats an integer, so it eventually leaves each finite set. Let $s_0>3$ be the minimum of a sufficiently late tail. After the occurrence of $s_i$, take the minimum of the remaining tail as $s_{i+1}$. The $s_i$ are strictly increasing and each survives forever. Each is odd, since an even value immediately halves. None is $3\bmod4$, since $(s_i+3)/2<s_i$. From each $s_i\equiv1\bmod4$ the next three steps reach $(5s_i+3)/4$. Therefore
\[
 s_{i+1}+3\le\tfrac54(s_i+3),\qquad s_i+3\le(5/4)^i(s_0+3).
\]
Every $i\ge0$ with $(5/4)^i(s_0+3)\le X+3$ supplies a distinct permanent survivor below or at $X$. Counting these indices proves both assertions.
\end{proof}
The census gives Haar survival exponent $1-h$, with $h=0.580684\ldots$ for this map. At $b_m=\lceil 2.4m\rceil$ its predictor $2^m P(b_m)=2^{(1-2.4(1-h))m+o(m)}$ tends to zero. Thus an upper bound of $3m$ actual survivors at that one depth for infinitely many $m$ would settle eventual periodicity. That bound remains open. The argument allows undiscovered cycles.

There is also a local test. On a divergent orbit take the last $s_i<X$. Since $(5/4)^3<2$, the next three tail minima all lie in $[X,2X)$ for sufficiently large $X$. Thus at most two survivors in $[2^m,2^{m+1})$ at the chosen depth for infinitely many $m$ would suffice as well. A single finite octave cannot establish this condition.
\newpage
\section{An explicit divergent orbit in negative mean drift}
Write $\F_4=\{0,1,w,w^2\}$ with $w^2=w+1$. The census map Q4b on $\F_4[t]$ divides by $t$ once per step. For $g=1+t^3$ it is
\[
 Q(f)=\begin{cases}
 f/t & f(0)=0,\\ g(f-1)/t & f(0)=1,\\ (f-w)/t+1 & f(0)=w,\\ (f-w^2)/t & f(0)=w^2 .
 \end{cases}
\]
At large degree the four degree changes are $-1,2,-1,-1$, giving mean $-1/4$ under uniform constant terms. Earlier census work disproved the polynomial analogue of E by counting periodic survivors. The existence of a divergent Q4b orbit was left open. The binary subring $\F_2[t]$ is invariant. On it we can transfer the four-phase binary orbit from the census note \emph{Provably divergent orbits}, D1. The following calculation proves the transfer directly.
\begin{theorem}
For $u=(1+t^3)(1+t+t^2+t^4+t^6)$ and every $j\ge0$,
\[
 \deg Q^j(u)\ge 9+j/2,\qquad
 \deg Q^{18(2^k-1)}(u)=9\cdot2^k\quad(k\ge0).
\]
In particular Q4b has a divergent orbit. For every even degree $D\ge18$ there is a polynomial on this orbit whose degree never subsequently falls below $D$.
\end{theorem}
\begin{proof}
Put $d=g+t$, $P(x)=1+x^6+x^9$, and $u_N=gP(t^N)/d$ for $N=2^k$.
Since $P(t)=d(1+t+t^2+t^4+t^6)$ and $P(t^N)=P(t)^N$, each $u_N$ is a polynomial of degree $9N$ and $u_1=u$.
Use the coordinate $z=df/g$. Its constant term equals that of $f$. On the binary subring the two updates in this coordinate are $z/t$ at constant term zero and $1+g(z+1)/t$ at constant term one.

Starting at $z=P(t^N)$, set $Y=t^{3N}$. The table gives the state at the end of each block. A block of zeros of length $r$ divides $z$ by $t^r$. A block of ones gives $1+g^r(z+1)/t^r$. Each length is the exact $t$-adic order of the entering $z$ or $z+1$, respectively, since $g(0)=1$. Frobenius gives $g^N=1+Y$.
\[
\begin{array}{ccl}
\toprule
 f(0)&\text{steps}&\text{ending }z\\\midrule
 1&6N&Y+Y^2+\cdots+Y^7\\
 0&3N&1+Y+\cdots+Y^6\\
 1&3N&Y^2+Y^6+Y^8\\
 0&6N&1+Y^4+Y^6=P(t^{2N})\\\bottomrule
\end{array}
\]
Consequently $Q^{18N}(u_N)=u_{2N}$. The boundary degrees are $9N$, $21N$, $18N$, $24N$ and $18N$. They vary linearly between boundaries, and lie on or above the line of slope $1/2$ joining the endpoints. Summing the block lengths from $N=1$ proves the asserted bound. In the first rising block every degree $9N+2i\le18N$ survives forever, since both later troughs have degree $18N$ and all later blocks start at or above it. The intervals $[9N,18N]$ for $N=2,4,8,\ldots$ cover every even degree from 18 on.
\end{proof}
\newpage
\section{Which negative integers cross zero for \texttt{split\_9}?}
On $\mathbb Z$ let $S(n)=n/2$ at even $n$, $S(n)=3n+9$ at $n\equiv1\pmod4$, and $S(n)=n+9$ at $n\equiv3\pmod4$. Eventual periodicity was already known. The negative basin for $d=9$ was recorded in the census as a finite sample, with 2,659 of the first 3,000 negative starts reaching zero. An exact count comes from stopping at the first change of sign or the first multiple of three.
\begin{theorem}
Let $\mathcal B$ consist of all integers obtained from either $1$ or $5$ by a finite word, possibly empty, in $L(u)=2u$ and $R(u)=4u+9$. For $u>0$, the orbit of $-u$ becomes positive if and only if $u\in\mathcal B$. In that case its eventual cycle is $(3,12,6)$. Every other negative start reaches zero.

If $B(X)=\#(\mathcal B\cap[1,X])$ and $\varphi=(1+\sqrt5)/2$, then
\[
 B(X)=\Theta\bigl(X^{\log_2\varphi}\bigr).
\]
\end{theorem}
\begin{proof}
On magnitudes $u=-n$ use one halving at even inputs and take the two forced halvings after each odd update. Until the first crossing this gives
\[
 u\longmapsto
 \begin{cases}
 u/2&u\text{ even},\\
 (u-9)/4&u\equiv1\pmod4,\\
 (3u-9)/4&u\equiv3\pmod4 .
 \end{cases}
\]
Every positive magnitude strictly decreases. A crossing into positive $n$ occurs precisely from $u=1$ or $5$, leading in the raw map to $8$ or $4$. Both reach $1\to12\to6\to3\to12$. Multiples of three are invariant and cannot reach either crossing seed. The third branch enters those multiples, so every pre-crossing path consists solely of the first two branches. Their positive inverse branches are exactly $L$ and $R$. Conversely each such inverse word follows its prescribed path, with every earlier magnitude positive. Distinct words and seeds have distinct outputs, since even outputs have unique parent under $L$, odd nonseeds in the tree have unique parent under $R$, and neither seed has a parent.

Give $L$ weight one and $R$ weight two. There are $F_{k+1}$ words of weight $k$, where $F_1=F_2=1$, by the recurrence for compositions into ones and twos. Starting at $s\in\{1,5\}$, a word of weight $k$ ends between $2^k$ and $8\cdot2^k$. Indeed, after dividing the output by $2^k$, the initial value is $s$ and each $R$ adds $9/2^j$, where $j$ is its cumulative weight. The $i$th such $j$ is at least $2i$, so the total addition is at most $3$. Summing the Fibonacci counts through $\lfloor\log_2 X\rfloor-3$ and through $\lfloor\log_2 X\rfloor$ gives the lower and upper bounds.
\end{proof}
Thus almost every negative start reaches zero, although infinitely many enter the positive cycle. Here the counted event is a basin crossing. In Theorem~\ref{thm:records} the counted event is survival above the starting height, and proving the required upper bound is the remaining integer convergence problem.

\smallskip\noindent\small\raggedright
\emph{Sources.} Collatz census (paths relative to \texttt{census/}).\\
\texttt{round3/deep/split5\_3\_deep.md} (Haar exponent),
\texttt{round3/polynomial/divergence.md} (D1), \texttt{round5/funcfield} (Q4b) and
\texttt{round3/domains/split\_family.md} (Theorem S2).\\
For prior work on polynomial stopping times see D. Nichols, \emph{Analogues of the $3x+1$ problem in polynomial rings of characteristic 2}, \href{https://arxiv.org/abs/1610.02545}{arXiv:1610.02545}.
\end{document}
